% from ega1-2-fr.tex
\oldpage[I]{99}

Par passage aux quotients, l'application
$\theta_x^\sharp:\mathscr{O}_{\psi(x)}\to\mathscr{O}_x$ donne donc un
monomorphisme $\theta^x:k(\psi(x))\to k(x)$, qui permet de considérer $k(x)$
comme une \emph{extension} du corps $k(\psi(x))$.

\begin{env}[2.2.2]
\label{I.2.2.2-fr}
Le composé $(\psi'',\theta'')$ de deux morphismes de préschémas
$(\psi,\theta)$ et $(\psi',\theta')$ est encore un morphisme de préschémas,
comme il résulte de la formule
${\theta''}^\sharp=\theta^\sharp\circ\psi^*({\theta'}^\sharp)$
(\hyperref[0.3.5.5-fr]{0, 3.5.5}). On en conclut que les préschémas forment
une \emph{catégorie}~; suivant la notation générale, on écrira
$\operatorname{Hom}(X,Y)$ l'ensemble des morphismes d'un préschéma $X$ dans
un préschéma $Y$.
\end{env}

\begin{example}[2.2.3]
\label{I.2.2.3-fr}
Si $U$ est une partie ouverte de $X$, l'injection canonique
(\hyperref[0.4.1.2-fr]{0, 4.1.2}) du préschéma induit
$(U,\mathscr{O}_X|U)$ dans $(X,\mathscr{O}_X)$ est un morphisme de
préschémas~; c'est d'ailleurs un \emph{monomorphisme} d'espaces annelés (et
\emph{a fortiori} un monomorphisme de préschémas), comme il résulte aussitôt
de (\hyperref[0.4.1.1-fr]{0, 4.1.1}).
\end{example}

\begin{proposition}[2.2.4]
\label{I.2.2.4-fr}
Soient $(X,\mathscr{O}_X)$ un préschéma, $(S,\mathscr{O}_S)$ un schéma affine
d'anneau $A$. Il existe une correspondance biunivoque canonique entre les
morphismes du préschéma $(X,\mathscr{O}_X)$ dans le préschéma
$(S,\mathscr{O}_S)$ et les homomorphismes de $A$ dans l'anneau
$\Gamma(X,\mathscr{O}_X)$.
\end{proposition}

Notons d'abord que, si $(X,\mathscr{O}_X)$ et $(Y,\mathscr{O}_Y)$ sont deux
espaces annelés quelconques, un morphisme $(\psi,\theta)$ de
$(X,\mathscr{O}_X)$ dans $(Y,\mathscr{O}_Y)$ définit canoniquement un
homomorphisme d'anneaux
$\Gamma(\theta):\Gamma(Y,\mathscr{O}_Y)\to
\Gamma(Y,\psi_*(\mathscr{O}_X))=\Gamma(X,\mathscr{O}_X)$. Dans le cas
considéré, tout revient à voir qu'un homomorphisme quelconque
$\varphi:A\to\Gamma(X,\mathscr{O}_X)$ est de la forme $\Gamma(\theta)$ pour
un $\theta$ et un seul. Or, il y a par hypothèse un recouvrement $(V_\alpha)$
de $X$ par des ouverts affines~; par composition de $\varphi$ avec
l'homomorphisme de restriction
$\Gamma(X,\mathscr{O}_X)\to
\Gamma(V_\alpha,\mathscr{O}_X|V_\alpha)$, on obtient un homomorphisme
$\varphi_\alpha:A\to\Gamma(V_\alpha,\mathscr{O}_X|V_\alpha)$, qui correspond
à un morphisme unique $(\psi_\alpha,\theta_\alpha)$ du préschéma
$(V_\alpha,\mathscr{O}_X|V_\alpha)$ dans $(S,\mathscr{O}_S)$, en vertu de
(\hyperref[I.1.7.3-fr]{1.7.3}). En outre, pour tout couple d'indices
$(\alpha,\beta)$, tout point de $V_\alpha\cap V_\beta$ admet un voisinage
ouvert affine $W$ contenu dans $V_\alpha\cap V_\beta$
(\hyperref[I.2.1.3-fr]{2.1.3})~; il est clair qu'en composant $\varphi_\alpha$
et $\varphi_\beta$ avec les homomorphismes de restriction à $W$, on obtient
le même homomorphisme
$\Gamma(S,\mathscr{O}_S)\to\Gamma(W,\mathscr{O}_X|W)$, donc, en vertu des
relations $(\theta_\alpha^\sharp)_x=(\varphi_\alpha)_x$ pour tout
$x\in V_\alpha$ et tout $\alpha$ (\hyperref[I.1.6.1-fr]{1.6.1}), les
restrictions à $W$ des morphismes $(\psi_\alpha,\theta_\alpha)$ et
$(\psi_\beta,\theta_\beta)$ coïncident. On en conclut qu'il y a un morphisme
d'espaces annelés
$(\psi,\theta):(X,\mathscr{O}_X)\to(S,\mathscr{O}_S)$ et un seul dont la
restriction à chaque $V_\alpha$ est $(\psi_\alpha,\theta_\alpha)$, et il est
clair que ce morphisme est un morphisme de préschémas et est tel que
$\Gamma(\theta)=\varphi$.

Soit $u:A\to\Gamma(X,\mathscr{O}_X)$ un homomorphisme d'anneaux, et soit
$v=(\psi,\theta)$ le morphisme $(X,\mathscr{O}_X)\to(S,\mathscr{O}_S)$
correspondant. Pour tout $f\in A$, on a
\begin{equation}
  \psi^{-1}(D(f))=X_{u(f)}
  \tag{2.2.4.1}\label{I.2.2.4.1-fr}
\end{equation}
avec les notations de (\hyperref[0.5.5.2-fr]{0, 5.5.2}) relatives au faisceau
localement libre $\mathscr{O}_X$. Il suffit en effet de vérifier cette formule
lorsque $X$ lui-même est affine, et alors elle n'est autre que
(\hyperref[I.1.2.2.2-fr]{1.2.2.2}).

\begin{proposition}[2.2.5]
\label{I.2.2.5-fr}
Sous les hypothèses de (\hyperref[I.2.2.4-fr]{2.2.4}), soient
$\varphi:A\to\Gamma(X,\mathscr{O}_X)$ un homomorphisme d'anneau,
$f:(X,\mathscr{O}_X)\to(S,\mathscr{O}_S)$ le morphisme de préschémas
correspondant, $\mathscr{G}$ (resp. $\mathscr{F}$) un
$\mathscr{O}_X$-Module (resp. un $\mathscr{O}_S$-Module) quasi-cohérent et
soit $M=\Gamma(S,\mathscr{F})$. Il existe alors une
